% =============================================================================
% Template: Projeto de Pesquisa — UnB
%
% Capa no padrão institucional clássico (Modelo B).
%
% Compilação local (a partir da raiz do repositório):
%   TEXINPUTS=./src: latexmk -pdf templates/unb-projeto-pesquisa/main.tex
% =============================================================================

\documentclass{abnt-unb-projeto-pesquisa}

\usepackage{abnt-citacoes}   % carrega abnt-refs automaticamente

\addbibresource{refs.bib}

% --- Metadados (preencha com seus dados) ---
\titulo{Título do Projeto de Pesquisa}
\subtitulo{subtítulo, se houver}
\autor{Nome Completo do Autor}
\orientador{Prof. Dr. Nome do Orientador}
\faculdade{Faculdade de Comunicação}
\departamento{Departamento de Comunicação}
\curso{Comunicação Social}
\unblogo{unb_logo}
\ano{2026}
\natureza{Projeto de pesquisa apresentado à Faculdade de Comunicação
  da Universidade de Brasília.}
\programa{Graduação em Comunicação Social}

% =============================================================================
\begin{document}

% --- Pré-textuais ---
\imprimircapa
\imprimirfolhaderosto
\imprimirsumario

% --- Textuais ---
\chapter{Introdução}

Contextualização e justificativa do problema de pesquisa \cite{eco1977}.

\chapter{Objetivos}

\section{Objetivo geral}

Descrever o objetivo geral.

\section{Objetivos específicos}

\begin{itemize}
  \item Objetivo específico 1
  \item Objetivo específico 2
  \item Objetivo específico 3
\end{itemize}

\chapter{Revisão de Literatura}

Revisão da literatura pertinente \cite{medeiros2019}.

\chapter{Metodologia}

Descrição dos métodos e procedimentos.

\chapter{Cronograma}

Cronograma previsto para a execução da pesquisa.

% --- Referências ---
\printbibliography

\end{document}
